\documentclass[10pt,a4paper]{amsart}
\usepackage{amsmath,amssymb,amsthm,graphicx}
% Retained minimal preamble for the Stacks Project chapter varieties.tex.
% Verbatim excerpt of the chapter preamble: Xy-pic setup (native preamble lines 14-18), the native theorem
% declarations with the mathematical macros (native preamble lines 154-217), and the project own no-op
% definitions of the reference/slogan/history display environments (native wrapper.tex lines 1-3) that the
% counted statement uses for its slogan. Copied mechanically; nothing added, removed or redefined.
% For commutative diagrams we use Xy-pic
\usepackage[all]{xy}

% We use 2cell for 2-commutative diagrams.
\xyoption{2cell}
% Theorem environments.
%
\theoremstyle{plain}
\newtheorem{theorem}[subsection]{Theorem}
\newtheorem{proposition}[subsection]{Proposition}
\newtheorem{lemma}[subsection]{Lemma}

\theoremstyle{definition}
\newtheorem{definition}[subsection]{Definition}
\newtheorem{example}[subsection]{Example}
\newtheorem{exercise}[subsection]{Exercise}
\newtheorem{situation}[subsection]{Situation}

\theoremstyle{remark}
\newtheorem{remark}[subsection]{Remark}
\newtheorem{remarks}[subsection]{Remarks}

\numberwithin{equation}{subsection}

% Macros
%
\def\lim{\mathop{\mathrm{lim}}\nolimits}
\def\colim{\mathop{\mathrm{colim}}\nolimits}
\def\Spec{\mathop{\mathrm{Spec}}}
\def\Hom{\mathop{\mathrm{Hom}}\nolimits}
\def\Ext{\mathop{\mathrm{Ext}}\nolimits}
\def\SheafHom{\mathop{\mathcal{H}\!\mathit{om}}\nolimits}
\def\SheafExt{\mathop{\mathcal{E}\!\mathit{xt}}\nolimits}
\def\Sch{\mathit{Sch}}
\def\Mor{\mathop{\mathrm{Mor}}\nolimits}
\def\Ob{\mathop{\mathrm{Ob}}\nolimits}
\def\Sh{\mathop{\mathit{Sh}}\nolimits}
\def\NL{\mathop{N\!L}\nolimits}
\def\CH{\mathop{\mathrm{CH}}\nolimits}
\def\proetale{{pro\text{-}\acute{e}tale}}
\def\etale{{\acute{e}tale}}
\def\QCoh{\mathit{QCoh}}
\def\Ker{\mathop{\mathrm{Ker}}}
\def\Im{\mathop{\mathrm{Im}}}
\def\Coker{\mathop{\mathrm{Coker}}}
\def\Coim{\mathop{\mathrm{Coim}}}

% Boxtimes
%
\DeclareMathSymbol{\boxtimes}{\mathbin}{AMSa}{"02}

%
% Macros for moduli stacks/spaces
%
\def\QCohstack{\mathcal{QC}\!\mathit{oh}}
\def\Cohstack{\mathcal{C}\!\mathit{oh}}
\def\Spacesstack{\mathcal{S}\!\mathit{paces}}
\def\Quotfunctor{\mathrm{Quot}}
\def\Hilbfunctor{\mathrm{Hilb}}
\def\Curvesstack{\mathcal{C}\!\mathit{urves}}
\def\Polarizedstack{\mathcal{P}\!\mathit{olarized}}
\def\Complexesstack{\mathcal{C}\!\mathit{omplexes}}
% \Pic is the operator that assigns to X its picard group, usage \Pic(X)
% \Picardstack_{X/B} denotes the Picard stack of X over B
% \Picardfunctor_{X/B} denotes the Picard functor of X over B
\def\Pic{\mathop{\mathrm{Pic}}\nolimits}
\def\Picardstack{\mathcal{P}\!\mathit{ic}}
\def\Picardfunctor{\mathrm{Pic}}
\def\Deformationcategory{\mathcal{D}\!\mathit{ef}}
\newenvironment{reference}{}{}
\newenvironment{slogan}{}{}
\newenvironment{history}{}{}
\usepackage{cleveref}
\crefname{originallemma}{original lemma}{original lemmas}
\Crefname{originallemma}{Original lemma}{Original lemmas}
\makeatletter
\newlabel{properties-lemma-ring-affine-open-injective-local-ring}{{0BX3}{source}{}{sourcepointer-1}{}}
\newlabel{properties-lemma-ring-affine-open-injective-local-ring@cref}{{[originallemma][0][]0BX3}{[1][0][]source}{}{}{}}
\makeatother
\crefname{originallemma}{original lemma}{original lemmas}
\Crefname{originallemma}{Original lemma}{Original lemmas}
\makeatletter
\newlabel{lemma-geometrically-reduced-at-point}{{035W}{source}{}{sourcepointer-2}{}}
\newlabel{lemma-geometrically-reduced-at-point@cref}{{[originallemma][0][]035W}{[1][0][]source}{}{}{}}
\makeatother
\crefname{originallemma}{original lemma}{original lemmas}
\Crefname{originallemma}{Original lemma}{Original lemmas}
\makeatletter
\newlabel{algebra-lemma-subalgebra-separable}{{030T}{source}{}{sourcepointer-3}{}}
\newlabel{algebra-lemma-subalgebra-separable@cref}{{[originallemma][0][]030T}{[1][0][]source}{}{}{}}
\makeatother
\title{Geometric reducedness at a point of a locally Noetherian scheme spreads to an affine open neighbourhood}
\author{The Stacks Project Authors}
\date{stacks-project:varieties.tex:local-snapshot-20261011; GNU FDL 1.2 or later}
\begin{document}
\maketitle
\noindent\textit{Editorial note.} All delivered mathematical text is the author's own, copied from the pinned Stacks Project chapter file varieties.tex; no statement or proof step was rewritten, reordered or removed.\par\medskip
\noindent\textit{Editorial note.} The preamble is the chapter's own retained minimal preamble (sources/stacks-varieties/retained-preamble.tex), a verbatim excerpt of the native Xy-pic setup, of the native theorem declarations with mathematical macros, and of the project's own no-op definitions of the reference display environments; no macro was invented or redefined.\par\medskip
\noindent\textit{Editorial note.} The same-chapter reference is located at its \textbackslash{}label line inside the pinned chapter file; the two cross-chapter pointers are located in the project's own tag list staged in the snapshot (sources/stacks-chow/tags.txt), at the verbatim lines carrying their canonical tags.\par\medskip
\noindent\textit{Editorial note.} Overlap pre-check: this label occurs in no other staged chapter (run with agent-authoring/stk/card\_plan.py varieties before opening).\par\medskip
\noindent\textit{Allowed standard prerequisites (editorial).}\begin{itemize}
\item An affine open neighbourhood whose ring injects into the local ring. Properties, Lemma labelled properties-lemma-ring-affine-open-injective-local-ring (canonical Stacks tag 0BX3); invoked at proof lines 610-611. Supplies the affine open U with R = O\_X(U) -> O\_\{X,x\} injective, which is where local Noetherianity of X is used.
\item Geometric reducedness at a point in local-ring terms. Same chapter, Lemma labelled lemma-geometrically-reduced-at-point (source line 354, Stacks tag 035W); invoked at proof lines 614-615. Supplies the translation of the hypothesis into geometric reducedness of the local ring O\_\{X,x\} over k.
\item A subalgebra of a geometrically reduced algebra over a field is geometrically reduced. Algebra, Lemma labelled algebra-lemma-subalgebra-separable (canonical Stacks tag 030T); invoked at proof lines 617-618. Supplies that R, being a subalgebra of the geometrically reduced ring O\_\{X,x\}, is geometrically reduced, whence U is.
\end{itemize}\medskip
\ifdefined\hypertarget\hypertarget{sourcepointer-1}{}\fi
\noindent\textit{Source reference.} Original lemma 0BX3: sources/stacks-chow/tags.txt, source lines 14966--14966. Cross-chapter standard prerequisite invoked once at proof lines 610-611 to obtain the affine open U with an injective map into the local ring; the rest of the argument is retained in the counted span.\par\medskip
\ifdefined\hypertarget\hypertarget{sourcepointer-2}{}\fi
\noindent\textit{Source reference.} Original lemma 035W: sources/stacks-varieties/pinned-original-source.tex, source lines 354--354. Same-chapter pointer invoked at proof lines 614-615 for the local-ring characterisation of the hypothesis; the neighbourhood conclusion is retained in the counted span.\par\medskip
\ifdefined\hypertarget\hypertarget{sourcepointer-3}{}\fi
\noindent\textit{Source reference.} Original lemma 030T: sources/stacks-chow/tags.txt, source lines 3960--3960. Cross-chapter standard prerequisite invoked once at proof lines 617-618 for the subalgebra criterion; the passage from R to U is retained in the counted span.\par\medskip
\begin{lemma}
\label{lemma-Noetherian-geometrically-reduced-at-point}
Let $k$ be a field.
Let $X$ be a scheme over $k$.
Let $x \in X$.
Assume $X$ locally Noetherian and geometrically reduced at $x$.
Then there exists an open neighbourhood $U \subset X$ of $x$
which is geometrically reduced over $k$.
\end{lemma}

\begin{proof}
Assume $X$ locally Noetherian and geometrically reduced at $x$.
By Properties, Lemma
\ref{properties-lemma-ring-affine-open-injective-local-ring}
we can find an affine open neighbourhood $U \subset X$ of $x$ such that
$R = \mathcal{O}_X(U) \to \mathcal{O}_{X, x}$
is injective. By
Lemma \ref{lemma-geometrically-reduced-at-point} the assumption
means that $\mathcal{O}_{X, x}$ is geometrically reduced over $k$.
By Algebra, Lemma \ref{algebra-lemma-subalgebra-separable}
this implies that $R$ is geometrically reduced over $k$, which
in turn implies that $U$ is geometrically reduced.
\end{proof}

\end{document}
